\documentclass[11pt,a4paper]{article}

% --- Packages ---
\usepackage[margin=.75in]{geometry}
\usepackage[utf8]{inputenc}
\usepackage[T1]{fontenc}
\usepackage{lmodern} % Fix for font not found error
\usepackage{titlesec}
\usepackage{enumitem}
\usepackage{hyperref}
\usepackage{xcolor}
\usepackage{fancyhdr}
\usepackage{comment}
\usepackage{cleveref}
\usepackage[bottom]{footmisc}
\usepackage{graphicx}
\usepackage{longtable}
\usepackage{booktabs}
\usepackage{array}
\usepackage{calc}
\usepackage{etoolbox}

% Standard Pandoc setup
\providecommand{\tightlist}{%
  \setlength{\itemsep}{0pt}\setlength{\parskip}{0pt}}

\crefformat{section}{\S#2#1#3}
\crefformat{subsection}{\S#2#1#3}
\crefformat{subsubsection}{\S#2#1#3}
\crefformat{figure}{#2Figure~#1#3}
\crefformat{footnote}{#2\footnotemark[#1]#3}

\linespread{0.9} % Riduce lo spazio tra le linee

% --- Colors ---
\definecolor{primary}{RGB}{0, 0, 0}
\definecolor{links}{HTML}{0000EE}
\definecolor{blue}{HTML}{26428B}

\hypersetup{
    colorlinks=true,
    linkcolor=links,
    urlcolor=links,
    citecolor=links,
}

% --- Section Formatting ---
\titleformat{\section}
  {\normalfont\Large\bfseries\color{blue}}{\thesection}{1em}{}
\titleformat{\subsection}
{\normalfont\large\bfseries\color{blue}}{\thesubsection}{1em}{}
\titleformat{\subsubsection}
{\normalfont\bfseries\color{blue}}{\thesubsubsection}{1em}{}

% --- Header and Footer ---
\pagestyle{fancy}
\fancyhead[L]{\small Giuseppe Capasso}
\fancyhead[R]{\small firewall}
\fancyfoot[C]{\thepage}
\renewcommand{\headrulewidth}{0.4pt}

% --- Custom Title Command ---
\newcommand{\proposaltitle}[1]{
    \begin{center}
        \vspace*{0.1cm}
        {\LARGE \bfseries \color{blue} #1} \\[1.5ex]
        {\large \textbf{Giuseppe Capasso}} \\[1ex]
                \small{
            \vspace{0.1cm}
            \textbf{Materia:} firewall\\
        }
            \end{center}
    \vspace{0.3cm}
    \hrule height 0.5pt
    \vspace{0.5cm}
}

\begin{document}

\proposaltitle{Lab 3 - Configure a ZPF}

\section{Lab 3 - Configure a ZPF}\label{lab-3---configure-a-zpf}

We will work in pair tables. Each table will have a LAN and will have to
forward only the correct traffic class to the other one.

Topology

Even tables (2, 4, 6) will have to forward \textbf{FTP} traffic while
odd tables (1, 3, 5) will have to allow \textbf{HTTP} traffic.

\subsubsection{Running test services}\label{running-test-services}

You can use Docker to run test servers. For example, to run a web server
you can use:

\begin{verbatim}
docker run -p 80:80 nginx
\end{verbatim}

Open your browser on localhost. You should see the following:

Web server

While for an ftp server, the command will be (substitute the with a real
directory you want to share):

\begin{verbatim}
docker run --env FTP_PASS=123 --env FTP_USER=user  --publish 20-21:20-21/tcp --publish 40000-40009:40000-40009/tcp --volume <path-to-dir>:/home/user garethflowers/ftp-server
\end{verbatim}

You can check the configuration by running an ftp client. Open your file
explorer and put the following in your bar

Ftp

\subsection{Configuration steps}\label{configuration-steps}

\textbf{\emph{Assumptions}}: the topology does not need routing as the
networks are directly attached. In the initial state, all devices can be
reached through \texttt{ping} and hosts in T1 network can configure
switch in T6 network. Assume that the server has a static IP address
configured.

Some part of the Lab will require the two tables to interact and agree
on policies

\subsubsection{Zone configuration}\label{zone-configuration}

The two tables will interact to decide zone names. These names will be
used to create the \texttt{zone-pair} and to assign the interface with
the \texttt{zone-member\ securitu\ \textless{}zone-name\textgreater{}}
command.

Configure the zones in the router with the
\texttt{zone\ security\ \textless{}zone-name\textgreater{}}.

\subsubsection{Traffic type
identification}\label{traffic-type-identification}

Now each group will create a \textbf{class map} to identify the traffic
to forward.

\textbf{\emph{Remember}}: a default class action exists and drops every
packet by default.

The traffic type should have two elements: 1. identify source and
destination IP addresses: for this purpose you should create an
\textbf{extended} ACL allowing for IP protocol specifying as source the
LAN source and as destination the static server IP address; 2. allow
HTTP (or FTP traffic): use multiple \textbf{match} statements when
configuring the class map;

For example:

\begin{verbatim}
r0(config)#access-list 100 permit ip 192.168.1.0 0.0.0.255 host 192.168.6.50
r0(config)#class-map type inspect match-all t1-traffic
r0(config-cmap)#match access-group 100
r0(config-cmap)#match protocol http
r0(config-cmap)#exit
r0(config)#
\end{verbatim}

\subsubsection{Create a policy}\label{create-a-policy}

Now you should create a policy and attach the created class map to it.

\begin{verbatim}
R0(config)#policy-map type inspect t1-policy
R0(config-pmap)#class type inspect t1-traffic
R0(config-pmap-c)#inspect
R0(config-pmap-c)#
\end{verbatim}

\subsubsection{Create a zone-pair}\label{create-a-zone-pair}

Now you can create a pair and assign the policy to it. For example, we
will configure the policy from t1 to t6.

\begin{verbatim}
R0(config)#zone-pair security t1-2-t6 source <t1-zone> destination <t6-zone>
R0(config-sec-zone-pair)#service-policy type inspect t1-policy
R0(config-sec-zone-pair)#exit
\end{verbatim}

\subsubsection{Assign the zones to the
interfaces}\label{assign-the-zones-to-the-interfaces}

Finally you can assign the interfaces according to zones. For example,
if \texttt{\textless{}t1-zone\textgreater{}} is attached to interface
\texttt{g/0/0/0}, command should be:

\begin{verbatim}
R0(config)#interface g/0/0/0
R0(config)#zone-memeber security <t1-zone>
\end{verbatim}

\subsubsection{Verify}\label{verify}

Now you should be able to visit the web server on \textbf{192.168.6.50},
but you cannot access any devices with ICMP messages. Same should be for
FTP server for \textbf{192.168.1.50}.

\end{document}
